% Actual finite formal checks required by the source examples, not historical proof status.
\begin{annexitem}{annex:p13.r11.logical}{thm:p13-r11-logical-benchmarks}{Rain, valid inference patterns, and the normative disjunct}
\begin{proof}
For affirming the consequent, take \(r=0,w=1\). Both the material
implication \(r\to w\) and \(w\) are true, but \(r\) is false.
Thus those premises do not entail rain. This is a Boolean countermodel;
it asserts neither a weather frequency nor a diagnosis of a person.

For modus ponens, a valuation satisfying \(p\to q\) and \(p\)
cannot make \(q\) false, since \((p,q)=(1,0)\) is the false row
of material implication. Therefore it satisfies \(q\).
For modus tollens, a valuation satisfying \(p\to q\) and
\(\neg q\) has \(q=0\). If it had \(p=1\), that same false
row would contradict the implication premise. Hence \(p=0\),
which proves \(\neg p\). These arguments use the declared
truth functions and do not depend on empirical conformity to them.

For the is--ought example, every valuation with \(D=1\) makes
\(D\lor N\) true, for either value of \(N\). In particular
\(D=1,N=0\) satisfies both the descriptive premise and the
disjunctive conclusion while falsifying the normative atom. Thus
\(D\entails D\lor N\) but \(D\not\entails N\).
Treating \(N\) as an uninterpreted atom is deliberate: no deontic
axioms, ethical premises, or special norm-validity semantics were
assumed. The result refutes a blanket claim that no conclusion
containing normative vocabulary follows from a descriptive premise;
it does not refute every nontrivial is--ought principle or justify
the particular bank-socialization demand in the source's example.
\end{proof}
\end{annexitem}
